\abstract{
从文档构建的知识图谱可能有字段缺失和取值错误。我们利用源文本和字段规则，将图检查与修复结合。局部、图和来源检查组成评估概况。主路径先清理输入，再让语言模型生成图，最后按字段约束过滤三元组。我们还测试尺度感知路由，每次为一个候选编辑赋予权重并进行检查。共享提示和过滤前后的评分展示各步骤的变化。在包含 450 个注入缺陷的 225 张图上，诊断＋门控 修复了 98.00\% 的缺陷。在 225 张模型构建图上，三元组 F1 从 87.99\% 提高到 92.48\%。直接字段复制在这些记录上达到 100\%。在 60 张新收据的测试中，来源行索引使 F1 比 简单流程 高 4.76 个百分点（95\% 置信区间：2.50 至 7.26）。在 56 份 CUAD 合同上，索引修复的 F1 为 44.84\%，输入为 46.76\%；与索引重新抽取的 F1 差为 0.72 个百分点（95\% 置信区间：$-4.80$ 至 6.48）。两名评审检查 200 项参考差异和 200 项系统编辑，在 200 项差异中的 95 项发现输入错误。第三名评审解决分歧后，五种设置中的 199 项确定编辑有 134 项被接受，另有一项仍不确定。实验展示源文本、字段检查和输入图如何影响修复。代码、模型输出和评价规则已提供，供后续研究使用。
}
\keywords{知识图谱修复，源文本，字段检查，文档抽取}